% Generated. Do not edit.
% Projection scriptural-date-location/section, version 1: the Scriptural Date and
% Location page, generated whole from research/chronology.toml, the Scripture
% chronology corpus and the source library by tools/tpt document-types refresh.
% document: liturgy/roman-rite/1962/propers/temporal/f01-christ-the-king
\runninghead{Scriptural Date and Location}
\section{Scriptural Date and Location}
\label{sec:date-location}

{\small Each passage the Mass reads or sings from Scripture, in canonical order. The Context column names what is dated --- the event the passage narrates, a setting it remembers, the traditional attribution of its author, or the composition of its book --- and where it belongs; the Date column gives the answer of the Scripture chronology corpus for that subject, with every alternative it records. Texts of the Mass that are not drawn from Scripture have no row. The evidence for each date follows the table.\par}
\medskip
% Generated. Do not edit.
% Deterministic projection of research/chronology.toml.
% schema: 5
% document: liturgy/roman-rite/1962/propers/temporal/f01-christ-the-king
% generated-by: tools/tpt proper-chronology annotations
\newcommand{\chronologyannotationclaim}[6]{#6}
\newcommand{\chronologyannotationreach}[3]{}
\newcommand{\chronologyannotationgroup}[3]{#3}
\newcommand{\chronologyannotation}[1]{%
  \ifcsname triptychchronologyannotation@#1\endcsname
    \csname triptychchronologyannotation@#1\endcsname
  \else
    \PackageError{triptych}{No chronology annotation for #1}{}%
  \fi}
% introit: Introit
\expandafter\def\csname triptychchronologyannotation@introit\endcsname{%
  \chronologyannotationgroup{historical-setting}{research-pending}{\textbf{Historical setting} -- Occasion date unresolved.}%
  \space \chronologyannotationgroup{composition}{nonuniform}{\textbf{Composition} -- No single date applies across every cited locus.}%
}
% epistle: Epistle
\expandafter\def\csname triptychchronologyannotation@epistle\endcsname{%
  \chronologyannotationgroup{composition}{disputed}{\textbf{Composition} -- disputed: \chronologyannotationreach{Col.1.12}{Col}{true}\chronologyannotationreach{Col.1.13}{Col}{true}\chronologyannotationreach{Col.1.14}{Col}{true}\chronologyannotationreach{Col.1.15}{Col}{true}\chronologyannotationreach{Col.1.16}{Col}{true}\chronologyannotationreach{Col.1.17}{Col}{true}\chronologyannotationreach{Col.1.18}{Col}{true}\chronologyannotationreach{Col.1.19}{Col}{true}\chronologyannotationreach{Col.1.20}{Col}{true}\chronologyannotationclaim{composition.epistle-to-the-colossians}{composition}{catholic-traditional-v1}{disputed}{(Philemon; Colossians; Ephesians; Philippians), 61}{A.D. 61}; \chronologyannotationreach{Col.1.12}{Col}{true}\chronologyannotationreach{Col.1.13}{Col}{true}\chronologyannotationreach{Col.1.14}{Col}{true}\chronologyannotationreach{Col.1.15}{Col}{true}\chronologyannotationreach{Col.1.16}{Col}{true}\chronologyannotationreach{Col.1.17}{Col}{true}\chronologyannotationreach{Col.1.18}{Col}{true}\chronologyannotationreach{Col.1.19}{Col}{true}\chronologyannotationreach{Col.1.20}{Col}{true}\chronologyannotationclaim{composition.epistle-to-the-colossians}{composition}{catholic-traditional-v1}{disputed}{some time between A.D. 61 and 63}{A.D. 61--63}.}%
}
% gradual: Gradual
\expandafter\def\csname triptychchronologyannotation@gradual\endcsname{%
  \chronologyannotationgroup{traditional-attribution}{disputed}{\textbf{Traditional attribution} -- disputed: \chronologyannotationreach{Ps.71.11}{Ps.2.8 Ps.28.10-11 Ps.71.1 Ps.71.8 Ps.71.11}{false}\chronologyannotationreach{Ps.71.8}{Ps.2.8 Ps.28.10-11 Ps.71.1 Ps.71.8 Ps.71.11}{false}\chronologyannotationclaim{israel.monarchy.david-traditional-era}{traditional-attribution}{catholic-traditional-v1}{disputed}{reigned from 1055 to 1015 B.C.}{David (reign in the usual chronology), B.C. 1055--1015}.}%
  \space \chronologyannotationgroup{historical-setting}{research-pending}{\textbf{Historical setting} -- Occasion date unresolved.}%
  \space \chronologyannotationgroup{composition}{preferred}{\textbf{Composition}: \chronologyannotationreach{Ps.71.11}{Ps}{true}\chronologyannotationreach{Ps.71.8}{Ps}{true}\chronologyannotationclaim{critical.psalms.latest-composition-boundary}{composition}{catholic-critical-v1}{preferred}{before the Maccabean period, around 165 B.C.; no individual psalm can be dated securely}{Before c. 165 B.C.}}%
}
% alleluia: Alleluia
\expandafter\def\csname triptychchronologyannotation@alleluia\endcsname{%
  \chronologyannotationgroup{historical-setting}{preferred}{\textbf{Historical setting}: \chronologyannotationreach{Dan.7.14}{Dan}{true}\chronologyannotationclaim{israel.exile.seventy-years}{historical-setting}{catholic-traditional-v1}{preferred}{seventy years}{The seventy years of servitude to Babylon, ``seventy years'' (duration)}.}%
  \space \chronologyannotationgroup{composition}{preferred}{\textbf{Composition}: \chronologyannotationreach{Dan.7.14}{Dan}{true}\chronologyannotationclaim{composition.book-of-daniel}{composition}{catholic-traditional-v1}{preferred}{admits 570-536 B.C. as its date of composition}{B.C. 570--536}.}%
}
% gospel: Gospel
\expandafter\def\csname triptychchronologyannotation@gospel\endcsname{%
  \chronologyannotationgroup{narrated-event}{research-pending}{\textbf{Event} -- No narrated-event date in the chronology corpus.}%
  \space \chronologyannotationgroup{composition}{disputed}{\textbf{Composition} -- disputed: \chronologyannotationreach{John.18.33}{John}{true}\chronologyannotationreach{John.18.34}{John}{true}\chronologyannotationreach{John.18.35}{John}{true}\chronologyannotationreach{John.18.36}{John}{true}\chronologyannotationreach{John.18.37}{John}{true}\chronologyannotationclaim{composition.gospel-of-john}{composition}{catholic-traditional-v1}{disputed}{from the year 90 to 100 (approximately)}{c. A.D. 90--100}; \chronologyannotationreach{John.18.33}{John}{true}\chronologyannotationreach{John.18.34}{John}{true}\chronologyannotationreach{John.18.35}{John}{true}\chronologyannotationreach{John.18.36}{John}{true}\chronologyannotationreach{John.18.37}{John}{true}\chronologyannotationclaim{composition.gospel-of-john}{composition}{catholic-traditional-v1}{disputed}{96 or one of the succeeding years}{A.D. 96--100}.}%
}
% offertory: Offertory
\expandafter\def\csname triptychchronologyannotation@offertory\endcsname{%
  \chronologyannotationgroup{traditional-attribution}{disputed}{\textbf{Traditional attribution} -- disputed: \chronologyannotationreach{Ps.2.8}{Ps.2.8 Ps.28.10-11 Ps.71.1 Ps.71.8 Ps.71.11}{false}\chronologyannotationclaim{israel.monarchy.david-traditional-era}{traditional-attribution}{catholic-traditional-v1}{disputed}{reigned from 1055 to 1015 B.C.}{David (reign in the usual chronology), B.C. 1055--1015}.}%
  \space \chronologyannotationgroup{historical-setting}{research-pending}{\textbf{Historical setting} -- Occasion date unresolved.}%
  \space \chronologyannotationgroup{composition}{preferred}{\textbf{Composition}: \chronologyannotationreach{Ps.2.8}{Ps}{true}\chronologyannotationclaim{critical.psalms.latest-composition-boundary}{composition}{catholic-critical-v1}{preferred}{before the Maccabean period, around 165 B.C.; no individual psalm can be dated securely}{Before c. 165 B.C.}}%
}
% communion: Communion
\expandafter\def\csname triptychchronologyannotation@communion\endcsname{%
  \chronologyannotationgroup{traditional-attribution}{disputed}{\textbf{Traditional attribution} -- disputed: \chronologyannotationreach{Ps.28.10}{Ps.2.8 Ps.28.10-11 Ps.71.1 Ps.71.8 Ps.71.11}{false}\chronologyannotationreach{Ps.28.11}{Ps.2.8 Ps.28.10-11 Ps.71.1 Ps.71.8 Ps.71.11}{false}\chronologyannotationclaim{israel.monarchy.david-traditional-era}{traditional-attribution}{catholic-traditional-v1}{disputed}{reigned from 1055 to 1015 B.C.}{David (reign in the usual chronology), B.C. 1055--1015}.}%
  \space \chronologyannotationgroup{historical-setting}{research-pending}{\textbf{Historical setting} -- Occasion date unresolved.}%
  \space \chronologyannotationgroup{composition}{preferred}{\textbf{Composition}: \chronologyannotationreach{Ps.28.10}{Ps}{true}\chronologyannotationreach{Ps.28.11}{Ps}{true}\chronologyannotationclaim{critical.psalms.latest-composition-boundary}{composition}{catholic-critical-v1}{preferred}{before the Maccabean period, around 165 B.C.; no individual psalm can be dated securely}{Before c. 165 B.C.}}%
}
% presentation: orientation-v3
\newcommand{\chronologyorientationrow}[2]{%
  \ifcsname triptychchronologyorientation@#1\endcsname
    \csname triptychchronologyorientation@#1\endcsname{#2}%
  \else
    \PackageError{triptych}{No chronology orientation row for #1}{}%
  \fi}
\expandafter\def\csname triptychchronologyorientation@row-e63c90797a12c4ac\endcsname#1{%
Offertory: Ps. 2:8 (Vg./Heb.) & Traditional attribution: David's reign; Hebron, then Jerusalem & Disputed: \chronologyannotationreach{Ps.2.8}{Ps.2.8 Ps.28.10-11 Ps.71.1 Ps.71.8 Ps.71.11}{false}\chronologyannotationclaim{israel.monarchy.david-traditional-era}{traditional-attribution}{catholic-traditional-v1}{disputed}{reigned from 1055 to 1015 B.C.}{B.C.~1055--1015}\\*[.1em]%
}
\expandafter\def\csname triptychchronologyorientation@row-8d0be5ee401d1031\endcsname#1{%
 & Historical setting & Occasion date unresolved\\*[.1em]%
}
\expandafter\def\csname triptychchronologyorientation@row-db2fc2d9dd69c14b\endcsname#1{%
 & Composition: Psalms & \chronologyannotationreach{Ps.2.8}{Ps}{true}\chronologyannotationclaim{critical.psalms.latest-composition-boundary}{composition}{catholic-critical-v1}{preferred}{before the Maccabean period, around 165 B.C.; no individual psalm can be dated securely}{Before c.~165 B.C.}\\[.3em]%
}
\expandafter\def\csname triptychchronologyorientation@row-7a3884a227d7ed18\endcsname#1{%
Communion: Ps. 28:10--11 (Vg.; Heb. 29) & Traditional attribution: David's reign; Hebron, then Jerusalem & Disputed: \chronologyannotationreach{Ps.28.10}{Ps.2.8 Ps.28.10-11 Ps.71.1 Ps.71.8 Ps.71.11}{false}\chronologyannotationreach{Ps.28.11}{Ps.2.8 Ps.28.10-11 Ps.71.1 Ps.71.8 Ps.71.11}{false}\chronologyannotationclaim{israel.monarchy.david-traditional-era}{traditional-attribution}{catholic-traditional-v1}{disputed}{reigned from 1055 to 1015 B.C.}{B.C.~1055--1015}\\*[.1em]%
}
\expandafter\def\csname triptychchronologyorientation@row-40853f150c4bee85\endcsname#1{%
 & Historical setting & Occasion date unresolved\\*[.1em]%
}
\expandafter\def\csname triptychchronologyorientation@row-0d1337853d1d6186\endcsname#1{%
 & Composition: Psalms & \chronologyannotationreach{Ps.28.10}{Ps}{true}\chronologyannotationreach{Ps.28.11}{Ps}{true}\chronologyannotationclaim{critical.psalms.latest-composition-boundary}{composition}{catholic-critical-v1}{preferred}{before the Maccabean period, around 165 B.C.; no individual psalm can be dated securely}{Before c.~165 B.C.}\\[.3em]%
}
\expandafter\def\csname triptychchronologyorientation@row-0aa553af61f123ed\endcsname#1{%
Introit: Ps 71:1 (Vg.; Heb. 72) & Traditional attribution: David's reign; Hebron, then Jerusalem & Disputed: \chronologyannotationreach{Ps.71.1}{Ps.2.8 Ps.28.10-11 Ps.71.1 Ps.71.8 Ps.71.11}{false}\chronologyannotationclaim{israel.monarchy.david-traditional-era}{traditional-attribution}{catholic-traditional-v1}{disputed}{reigned from 1055 to 1015 B.C.}{B.C.~1055--1015}\\*[.1em]%
}
\expandafter\def\csname triptychchronologyorientation@row-26196dc4428c5d03\endcsname#1{%
 & Historical setting & Occasion date unresolved\\*[.1em]%
}
\expandafter\def\csname triptychchronologyorientation@row-9757e96e6962c0e3\endcsname#1{%
 & Composition: Psalms & \chronologyannotationreach{Ps.71.1}{Ps}{true}\chronologyannotationclaim{critical.psalms.latest-composition-boundary}{composition}{catholic-critical-v1}{preferred}{before the Maccabean period, around 165 B.C.; no individual psalm can be dated securely}{Before c.~165 B.C.}\\[.3em]%
}
\expandafter\def\csname triptychchronologyorientation@row-baecd1b713331c34\endcsname#1{%
Gradual: Ps. 71:8,11 (Vg.; Heb. 72) & Traditional attribution: David's reign; Hebron, then Jerusalem & Disputed: \chronologyannotationreach{Ps.71.11}{Ps.2.8 Ps.28.10-11 Ps.71.1 Ps.71.8 Ps.71.11}{false}\chronologyannotationreach{Ps.71.8}{Ps.2.8 Ps.28.10-11 Ps.71.1 Ps.71.8 Ps.71.11}{false}\chronologyannotationclaim{israel.monarchy.david-traditional-era}{traditional-attribution}{catholic-traditional-v1}{disputed}{reigned from 1055 to 1015 B.C.}{B.C.~1055--1015}\\*[.1em]%
}
\expandafter\def\csname triptychchronologyorientation@row-42915bce5b517b01\endcsname#1{%
 & Historical setting & Occasion date unresolved\\*[.1em]%
}
\expandafter\def\csname triptychchronologyorientation@row-0935edd74b354160\endcsname#1{%
 & Composition: Psalms & \chronologyannotationreach{Ps.71.11}{Ps}{true}\chronologyannotationreach{Ps.71.8}{Ps}{true}\chronologyannotationclaim{critical.psalms.latest-composition-boundary}{composition}{catholic-critical-v1}{preferred}{before the Maccabean period, around 165 B.C.; no individual psalm can be dated securely}{Before c.~165 B.C.}\\[.3em]%
}
\expandafter\def\csname triptychchronologyorientation@row-1c32568014a7f91a\endcsname#1{%
Alleluia: Dan. 7:14 & Historical setting: The seventy years of servitude to Babylon; Judah and Babylon & \chronologyannotationreach{Dan.7.14}{Dan}{true}\chronologyannotationclaim{israel.exile.seventy-years}{historical-setting}{catholic-traditional-v1}{preferred}{seventy years}{``seventy years'' (duration)}\\*[.1em]%
}
\expandafter\def\csname triptychchronologyorientation@row-a0c1bd99a7ed1394\endcsname#1{%
 & Composition: Daniel; Babylonian setting & \chronologyannotationreach{Dan.7.14}{Dan}{true}\chronologyannotationclaim{composition.book-of-daniel}{composition}{catholic-traditional-v1}{preferred}{admits 570-536 B.C. as its date of composition}{B.C.~570--536}\\[.3em]%
}
\expandafter\def\csname triptychchronologyorientation@row-87b866de388ce05f\endcsname#1{%
Gospel: John 18:33--37 & Narrated event: Jesus before Pilate in the governor's hall at Jerusalem & Date unresolved\\*[.1em]%
}
\expandafter\def\csname triptychchronologyorientation@row-ba70b39d98d421e9\endcsname#1{%
 & Composition: John; Ephesus traditionally & Disputed: \chronologyannotationreach{John.18.33}{John}{true}\chronologyannotationreach{John.18.34}{John}{true}\chronologyannotationreach{John.18.35}{John}{true}\chronologyannotationreach{John.18.36}{John}{true}\chronologyannotationreach{John.18.37}{John}{true}\chronologyannotationclaim{composition.gospel-of-john}{composition}{catholic-traditional-v1}{disputed}{from the year 90 to 100 (approximately)}{c.~A.D.~90--100}; \chronologyannotationreach{John.18.33}{John}{true}\chronologyannotationreach{John.18.34}{John}{true}\chronologyannotationreach{John.18.35}{John}{true}\chronologyannotationreach{John.18.36}{John}{true}\chronologyannotationreach{John.18.37}{John}{true}\chronologyannotationclaim{composition.gospel-of-john}{composition}{catholic-traditional-v1}{disputed}{96 or one of the succeeding years}{A.D.~96--100}\\[.3em]%
}
\expandafter\def\csname triptychchronologyorientation@row-c308e94d55cfa7ba\endcsname#1{%
Epistle: Col. 1:12--20 & Composition: Colossians; Rome traditionally; addressed to Colossae & Disputed: \chronologyannotationreach{Col.1.12}{Col}{true}\chronologyannotationreach{Col.1.13}{Col}{true}\chronologyannotationreach{Col.1.14}{Col}{true}\chronologyannotationreach{Col.1.15}{Col}{true}\chronologyannotationreach{Col.1.16}{Col}{true}\chronologyannotationreach{Col.1.17}{Col}{true}\chronologyannotationreach{Col.1.18}{Col}{true}\chronologyannotationreach{Col.1.19}{Col}{true}\chronologyannotationreach{Col.1.20}{Col}{true}\chronologyannotationclaim{composition.epistle-to-the-colossians}{composition}{catholic-traditional-v1}{disputed}{(Philemon; Colossians; Ephesians; Philippians), 61}{A.D.~61}; \chronologyannotationreach{Col.1.12}{Col}{true}\chronologyannotationreach{Col.1.13}{Col}{true}\chronologyannotationreach{Col.1.14}{Col}{true}\chronologyannotationreach{Col.1.15}{Col}{true}\chronologyannotationreach{Col.1.16}{Col}{true}\chronologyannotationreach{Col.1.17}{Col}{true}\chronologyannotationreach{Col.1.18}{Col}{true}\chronologyannotationreach{Col.1.19}{Col}{true}\chronologyannotationreach{Col.1.20}{Col}{true}\chronologyannotationclaim{composition.epistle-to-the-colossians}{composition}{catholic-traditional-v1}{disputed}{some time between A.D. 61 and 63}{A.D.~61--63}\\[.3em]%
}
\expandafter\def\csname triptychchronologyorientation@row-6298d6c79b112d82\endcsname#1{%
Introit: Apoc 1:6, Apoc 5:12 & Composition: Apocalypse; seer on Patmos; heavenly scenes are visions & Disputed: \chronologyannotationreach{Apoc.1.6}{Apoc}{true}\chronologyannotationreach{Apoc.5.12}{Apoc}{true}\chronologyannotationclaim{composition.apocalypse}{composition}{catholic-traditional-v1}{disputed}{the reign of Domitian ... A.D. 81-96}{A.D.~81--96}; \chronologyannotationreach{Apoc.1.6}{Apoc}{true}\chronologyannotationreach{Apoc.5.12}{Apoc}{true}\chronologyannotationclaim{composition.apocalypse}{composition}{catholic-traditional-v1}{disputed}{from the year 90 to 100 (approximately)}{c.~A.D.~90--100}\\[.3em]%
}

% Generated. Do not edit.
% Complete orientation-v3 table; wording, dates and order belong to the corpus.
\begin{chronologysharedtable}
\chronologyorientationrow{row-b65b2e47f275ce3d}{}
\chronologyorientationrow{row-592bd7ea4d8b4045}{}
\chronologyorientationrow{row-3d54b709c87afe27}{}
\midrule
\chronologyorientationrow{row-c18a401b8c1f247a}{}
\chronologyorientationrow{row-3a43a808252fefb6}{}
\chronologyorientationrow{row-0fa66f253aa1673a}{}
\midrule
\chronologyorientationrow{row-3fd7b7e5eade10d4}{}
\chronologyorientationrow{row-cc49b03a4d523790}{}
\chronologyorientationrow{row-b1e094f0018f10f0}{}
\midrule
\chronologyorientationrow{row-389bbaca1560ec2a}{}
\chronologyorientationrow{row-d843eca867531f5a}{}
\chronologyorientationrow{row-786e78ea24756854}{}
\midrule
\chronologyorientationrow{row-1c32568014a7f91a}{}
\chronologyorientationrow{row-a0c1bd99a7ed1394}{}
\midrule
\chronologyorientationrow{row-87b866de388ce05f}{}
\chronologyorientationrow{row-ba70b39d98d421e9}{}
\midrule
\chronologyorientationrow{row-c308e94d55cfa7ba}{}
\midrule
\chronologyorientationrow{row-6298d6c79b112d82}{}
\end{chronologysharedtable}
{\footnotesize Psalm numbering: Vg. = Vulgate; Heb. = Hebrew.\par}


\subsection{Evidence for each date}
\label{sec:date-location-evidence}
{\small For each passage, as the Scriptural Date and Location table names it: what is dated, how the corpus weighs each answer, the source's own words, and the sources cited. A subject the corpus does not date is reported as such and given no date.\par}
{\small
\par\smallskip\noindent\textit{Offertory: Ps. 2:8 (Vg./Heb.)}\par\nobreak
\begin{itemize}[leftmargin=1.2em,itemsep=1pt,topsep=1pt,parsep=0pt]
\item \textbf{Traditional attribution} (disputed). David's reign: B.C.~1055--1015 (``reigned from 1055 to 1015 B.C.''; The Catholic Encyclopedia, Volume 4 (New York, 1908), article 'David, King,' printed p. 642, opening paragraph). Place: Hebron, then Jerusalem (The Holy Bible, Douay-Rheims (Challoner revision), Project Gutenberg eBook 1581: 2 Kings 5:5).
\item \textbf{Historical setting}. No date: the chronology corpus has not yet dated this (research pending).
\item \textbf{Composition} (preferred). Psalms: Before c.~165 B.C. (``before the Maccabean period, around 165 B.C.; no individual psalm can be dated securely''; Catholic critical chronology; New American Bible, Revised Edition --- USCCB web state (2026), Psalms introduction).
\end{itemize}
\par\smallskip\noindent\textit{Communion: Ps. 28:10--11 (Vg.; Heb. 29)}\par\nobreak
\begin{itemize}[leftmargin=1.2em,itemsep=1pt,topsep=1pt,parsep=0pt]
\item \textbf{Traditional attribution}. As under Offertory: Ps. 2:8 (Vg./Heb.).
\item \textbf{Historical setting}. As under Offertory: Ps. 2:8 (Vg./Heb.).
\item \textbf{Composition}. As under Offertory: Ps. 2:8 (Vg./Heb.).
\end{itemize}
\par\smallskip\noindent\textit{Introit: Ps. 71:1 (Vg.; Heb. 72)}\par\nobreak
\begin{itemize}[leftmargin=1.2em,itemsep=1pt,topsep=1pt,parsep=0pt]
\item \textbf{Traditional attribution}. As under Offertory: Ps. 2:8 (Vg./Heb.).
\item \textbf{Historical setting}. As under Offertory: Ps. 2:8 (Vg./Heb.).
\item \textbf{Composition}. As under Offertory: Ps. 2:8 (Vg./Heb.).
\end{itemize}
\par\smallskip\noindent\textit{Gradual: Ps. 71:8,11 (Vg.; Heb. 72)}\par\nobreak
\begin{itemize}[leftmargin=1.2em,itemsep=1pt,topsep=1pt,parsep=0pt]
\item \textbf{Traditional attribution}. As under Offertory: Ps. 2:8 (Vg./Heb.).
\item \textbf{Historical setting}. As under Offertory: Ps. 2:8 (Vg./Heb.).
\item \textbf{Composition}. As under Offertory: Ps. 2:8 (Vg./Heb.).
\end{itemize}
\par\smallskip\noindent\textit{Alleluia: Dan. 7:14}\par\nobreak
\begin{itemize}[leftmargin=1.2em,itemsep=1pt,topsep=1pt,parsep=0pt]
\item \textbf{Historical setting} (preferred). The seventy years of servitude to Babylon: ``seventy years'' (duration) (The Holy Bible, Douay-Rheims (Challoner revision), Project Gutenberg eBook 1581: Jer. 25:11, 25:12). Place: Judah and Babylon (The Holy Bible, Douay-Rheims (Challoner revision), Project Gutenberg eBook 1581: Jer. 25:11, 25:12).
\item \textbf{Composition} (preferred). The Book of Daniel: B.C.~570--536 (``admits 570-536 B.C. as its date of composition''; The Catholic Encyclopedia, Volume 4 (New York, 1908)). Place: Babylonian setting (The Holy Bible, Douay-Rheims (Challoner revision), Project Gutenberg eBook 1581: Dan. 1:1, 1:6, 7:1).
\end{itemize}
\par\smallskip\noindent\textit{Gospel: John 18:33--37}\par\nobreak
\begin{itemize}[leftmargin=1.2em,itemsep=1pt,topsep=1pt,parsep=0pt]
\item \textbf{Narrated event}. Jesus before Pilate in the governor's hall at Jerusalem (The Holy Bible, Douay-Rheims (Challoner revision), Project Gutenberg eBook 1581: John 18:1, 18:28, 18:33, 18:37). No date: the chronology corpus has not yet dated this (research pending).
\item \textbf{Composition} (disputed). The Gospel of St John: c.~A.D.~90--100 (``from the year 90 to 100 (approximately)''; The Catholic Encyclopedia, Volume 14 (New York, 1912)); A.D.~96--100 (``96 or one of the succeeding years''; The Catholic Encyclopedia, Volume 8 (New York, 1910)). Place: Ephesus traditionally (The Catholic Encyclopedia, Volume 8 (New York, 1910)).
\end{itemize}
\par\smallskip\noindent\textit{Epistle: Col. 1:12--20}\par\nobreak
\begin{itemize}[leftmargin=1.2em,itemsep=1pt,topsep=1pt,parsep=0pt]
\item \textbf{Composition} (disputed). The Epistle to the Colossians: A.D.~61 (``(Philemon; Colossians; Ephesians; Philippians), 61''; The Catholic Encyclopedia, Volume 11 (New York, 1911)); A.D.~61--63 (``some time between A.D. 61 and 63''; The Catholic Encyclopedia, Volume 4 (New York, 1908)). Place: Rome traditionally; addressed to Colossae (The Catholic Encyclopedia, Volume 4 (New York, 1908); The Holy Bible, Douay-Rheims (Challoner revision), Project Gutenberg eBook 1581: Col. 1:2).
\end{itemize}
\par\smallskip\noindent\textit{Introit: Apoc. 1:6, Apoc. 5:12}\par\nobreak
\begin{itemize}[leftmargin=1.2em,itemsep=1pt,topsep=1pt,parsep=0pt]
\item \textbf{Composition} (disputed). The Apocalypse of St John: A.D.~81--96 (``the reign of Domitian ... A.D. 81-96''; The Catholic Encyclopedia, Volume 1 (New York, 1907)); c.~A.D.~90--100 (``from the year 90 to 100 (approximately)''; The Catholic Encyclopedia, Volume 14 (New York, 1912)). Place: seer on Patmos; heavenly scenes are visions (The Holy Bible, Douay-Rheims (Challoner revision), Project Gutenberg eBook 1581: Apoc. 1:9, 4:1, 7:9).
\end{itemize}
\par}
